\section{Discussion}
\label{sec:discussion}

The qualitative data lets us address the research question directly: in what ways does the paired-scenario game support participants' understanding of Nielsen's heuristics, and how do the surrounding design choices --- narration and light gamification --- affect that support? Four themes recur across the responses, and we discuss each in turn before stating the limitations of the study.

The first theme is that the paired bad/good contrast, combined with the interactive nature of the scenarios, is the mechanism participants most often name when explaining how the game supported their understanding (Section~\ref{sec:results:t2}). This pattern fits an experiential-learning reading of the activity: the player has a concrete encounter with a violating interface, then with a compliant one, and the reveal screen labels the abstract concept that explains the contrast. The sequence operationalises the move from concrete experience to abstract conceptualisation described by Kolb \cite{kolb1984experiential}, and the participants' own language --- $P_{30}$'s ``a more active method of learning than just reading, which is a more passive way'' --- mirrors the active-learning literature \cite{bonwell1991active,freeman2014active}. The contrast itself, named by participants as ``error and correct'' or ``worse versus better'', is what gives the abstract definition something concrete to attach to.

The second theme is that gamification operates as engagement scaffolding rather than as a separate learning mechanism. Participants who said gamification helped them did so by referring to focus, motivation, or the wish to finish the sequence --- not to understanding the heuristics. Several participants explicitly singled out the progress bar as helpful while reporting that the stars and the closing screen did not affect them ($P_2$, $P_5$, $P_{14}$, $P_{22}$, $P_{23}$). Two participants critiqued the gamification more structurally for not introducing a real challenge ($P_6$, $P_7$). The reading we take from this is that progress tracking carried most of the engagement value, and that the achievement layer is the most easily revised part of the design with little expected cost to the support reported in Section~\ref{sec:results:t2}.

The third theme is that the narration is useful for most participants and costly for some. The narration helped about a third of the sample focus and notice the heuristic in action (anchor: $P_8$). It distracted or frustrated a smaller group, with consistent complaints about the absence of a mute or volume control, the length of some tracks, and the speech style ($P_{18}$, $P_{20}$, $P_{22}$, $P_{10}$). The simplest design response is to add per-track playback controls and to make the audio non-automatic, so that each player can choose the modality that fits their attention style. This is not a question of removing narration, but of letting it be an opt-in rather than an opt-out feature.

The fourth theme is the recurring tension around whether the artefact is a game at all. Several participants ($P_2$, $P_7$, $P_{18}$, $P_{30}$) read it as a guided demonstration rather than as a game proper, on the grounds that there is no failure state and no real decision to be made. This critique is honest and consistent with the design choices reported in Section~\ref{sec:game}: the artefact intentionally avoids assessment, scoring, or branching. The participants' proposal --- that the player be allowed to identify the heuristic before the reveal, or to lose stars on mistakes ($P_2$, $P_8$) --- is the most actionable future-work direction we have. It would shift the artefact from a pure content-delivery device to one that also probes recognition, which is the cognitive step on which novices are reported to struggle \cite{nielsen1992finding,abulfaraj2020detailed}.

Several limitations bound the contribution of this experience report. The sample is small ($N = 30$) and self-selected; recruitment combined a classroom deployment with the research team's network, so the participants who chose to play and complete the questionnaire may have been more receptive to the activity than a random sample would have been. The sample is heavily concentrated in Computer Science and Software Engineering (24 of 30) and includes no participants from Education or from outside technology areas, which limits the generalisability of the findings to other audiences and to formal teacher-led settings. The study collected no pre- or post-test of heuristic knowledge, so the claims rest on self-reported support for understanding rather than on measured learning gains. Finally, the deployment was single-session, so retention is not assessed; a delayed re-evaluation would be a natural next step.
